\textbf{Proof.}
Let \(d=c-c'\).  Since both contrasts have zero sum, \(\sum_a d_a=0\).  Write
\[
 P_d=\{a:d_a>0\},\qquad N_d=\{a:d_a<0\}.
\]
Then
\[
 \sum_{a\in P_d}d_a=-\sum_{a\in N_d}d_a
 =\frac12\sum_{a\in\mathcal A_K}|d_a|=\eta(c,c').
\]
For every binary response type \(t\), selecting an arbitrary subset of the coordinates of a zero-sum vector gives
\[
 -\eta(c,c')\leq\sum_{a\in\mathcal A_K}d_at_a\leq\eta(c,c').
\]
Averaging this inequality over the \(n\) labeled units proves the uniform target bound
\[
 \sup_{z\in\mathcal T_K^n}|\tau_c(z)-\tau_{c'}(z)|
 \leq\eta(c,c').
\]
This calculation includes every boundary schedule.  In particular, the all-zero and all-one schedules give target difference zero.

Put \(I_c=[-L_c/2,L_c/2]\).  Fix any assignment law \(\mathcal D\) and any estimator \(f'\) admissible for contrast \(c'\), and define the estimator admissible for contrast \(c\) by
\[
 f=\operatorname{proj}_{I_c}(f').
\]
For a fixed schedule \(z\), the target \(\tau_c(z)\) belongs to \(I_c\), so projection cannot increase its pointwise distance.  The triangle inequality in \(L^2(\mathcal D)\), namely Minkowski's inequality, and the uniform target bound give
\[
 \begin{aligned}
 \sqrt{R_n(\mathcal D,f;z)}
 &\leq
 \left\|f'-\tau_c(z)\right\|_{L^2(\mathcal D)}\\
 &\leq
 \left\|f'-\tau_{c'}(z)\right\|_{L^2(\mathcal D)}
 +|\tau_{c'}(z)-\tau_c(z)|\\
 &\leq \sqrt{R_n(\mathcal D,f';z)}+\eta(c,c').
 \end{aligned}
\]
Taking the maximum over schedules, then the infimum over all procedures for \(c'\), yields
\[
 \sqrt{\rho_n(c)}\leq\sqrt{\rho_n(c')}+\eta(c,c').
\]
Interchanging \(c\) and \(c'\) gives the reverse one-sided inequality.  Combining them proves the displayed absolute-value bound.  Every projection and norm is well defined because both contrasts are nonzero, so both target intervals have positive finite radius. \(\square\)+